% poster.tex: an example A0 conference poster in Vellum (light), after "Vellum - conference posters".
% Palmer penguins (Gorman, Williams and Fraser 2014). The central figure and every number are made by
% examples/python/poster/make_poster_figures.py (fig/central.pdf, fig/numbers.tex). Build: python examples/latex/poster/build.py
% Sizes: headline 90 pt, authors 32 pt, column headings 36 pt, body 28 pt, stat 72 pt over 24 pt, footer 18 pt;
% the figure is placed at 1:1 (drawn at 2.6 x V.size("poster") with the type scaled by 2.6). Nothing is under 18 pt.
\documentclass{beamer}
\usepackage[orientation=portrait, size=a0, scale=1]{beamerposter}
\usefonttheme{professionalfonts}
\usepackage{fontspec}
\setmainfont{TeX Gyre Pagella}
\setsansfont{Lato}
\usepackage{tikz}
\usepackage{qrcode}
\usepackage{vellum-colors}
\vellumbeamer
\input{fig/numbers.tex}
\setbeamertemplate{footline}{}
\setbeamertemplate{headline}{}

% TikZ's align=left fixes word spaces at 1/3 em of the node's font; text set at another size inside a node takes
% its own font's space again (no stretch, so a ragged line stays ragged)
\newcommand\natspace{\spaceskip=\fontdimen2\font\relax\xspaceskip=0pt\relax}
\newcommand\headsize{\fontsize{90pt}{100pt}\selectfont\natspace}
\newcommand\authsize{\fontsize{32pt}{40pt}\selectfont\natspace}
\newcommand\colhead{\fontsize{36pt}{44pt}\selectfont\bfseries\natspace}
\newcommand\bodysize{\fontsize{28pt}{40pt}\selectfont\natspace}
\newcommand\statnum{\fontsize{72pt}{78pt}\selectfont\natspace}
\newcommand\statlabel{\fontsize{24pt}{31pt}\selectfont\natspace}
\newcommand\footsize{\fontsize{18pt}{24pt}\selectfont\natspace}
% one numbered point: the number hangs before the text, in the teal accent of the headings
\newcommand\point[2]{\par\noindent\hangindent=1.25em\hangafter=1\makebox[1.25em][l]{\color{vteal}\bfseries #1}#2\par\vskip 9mm}
% a column card at x, w mm wide: cream, no border, heading in 36 pt bold Lato, a numbered list at 28 pt, hung from the top
\newcommand\cardy{-695}
\newcommand\cardh{205}
\newcommand\card[4]{\fill[vCard, rounded corners=2mm] (#1, \cardy) rectangle ++(#2, -\cardh);
  \node[anchor=north west, inner sep=0pt, text width=#2mm-40mm, align=left, text=vInk] at (#1+20, \cardy-20)
    {{\colhead\color{vteal}#3\par}\vskip 12mm\bodysize\raggedright #4};}

\begin{document}
\sffamily\frenchspacing
\begin{frame}[plain]
\begin{tikzpicture}[remember picture, overlay, shift={(current page.north west)}, x=1mm, y=1mm]
  % ---- headline band: the claim, one line, and who made it
  \node[anchor=north west, inner sep=0pt, text=vInk] at (30, -30) {\rmfamily\headsize Beak shape, not body size, tells penguin species apart};
  \node[anchor=north west, inner sep=0pt, text=vSecondary] at (30, -78) {\authsize A.\ Kaminas\enspace·\enspace Vellum example poster\enspace·\enspace
    data: Gorman, Williams and Fraser (2014), Palmer Station LTER};
  % ---- the one figure, at its drawn size
  \node[anchor=north west, inner sep=0pt] at (30, -158) {\includegraphics{fig/central.pdf}};
  % ---- three columns in reading order: the question, what was done, what it means
  \card{30}{247}{The question}{%
    \point{1}{Three penguin species breed on the Palmer Archipelago: Adélie, Chinstrap and Gentoo.}
    \point{2}{Gentoo are heavy, \pMassGentoo\ g on average; Adélie and Chinstrap differ by only \pMassGap\ g.}
    \point{3}{Which measure tells all three apart: the size of the body, or the shape of the beak?}}
  \card{297}{247}{What was done}{%
    \point{1}{\pBirds\ birds measured at Palmer Station: beak length and depth, flipper length, body mass.}
    \point{2}{Each measure scaled to mean 0 and standard deviation 1.}
    \point{3}{Each bird left out in turn and assigned to the species with the nearest mean.}
    \point{4}{Once on beak length and depth, once on body mass alone.}}
  \card{564}{247}{What it means}{%
    \point{1}{Beak length and depth place \pAccBeak\,\% of birds in their own species; body mass places \pAccMass\,\%.}
    \point{2}{Adélie beaks are short and deep, Chinstrap long and deep, Gentoo long and shallow.}
    \point{3}{Size sets only Gentoo apart: flipper length alone places \pAccFlipper\,\%.}
    \point{4}{Beak shape is the measure that tells all three apart.}}
  % ---- the stat tile, where reading ends: the one number that matters, terracotta Pagella 72 pt, its label 24 pt
  \fill[vCard, rounded corners=2mm] (30, -930) rectangle (811, -1004);
  \node[anchor=west, inner sep=0pt] (stat) at (50, -967) {\rmfamily\statnum\color{vterracotta}\pAccBeak\,\%};
  \node[anchor=base west, inner sep=0pt] at ([xshift=14mm]stat.base east) {\statlabel\color{vSecondary}of birds placed in their own
    species by beak length and depth, against \pAccMass\,\% by body mass alone};
  % ---- footer band: contact, data, script and stamp at 18 pt muted; the QR code to the site
  \draw[vGrid, line width=0.8pt] (30, -1078) -- (811, -1078);
  \node[anchor=west, inner sep=0pt, text width=690mm, align=left, text=vMuted] at (30, -1132) {\footsize
    Contact: A.\ Kaminas\enspace·\enspace akaminas.github.io/vellum\par
    Data: Gorman, Williams and Fraser (2014), PLoS ONE 9(3): e90081; Palmer Station LTER; via the palmerpenguins R package and vega-datasets 2.8.0; CC0\par
    Accuracies: nearest species mean, leave-one-out, on \pBirds\ birds\enspace·\enspace make\_poster\_figures.py\enspace·\enspace poster.tex\enspace·\enspace git \pStamp\par};
  \node[anchor=south east, inner sep=0pt] at (811, -1159) {\color{vInk}\qrcode[height=54mm]{https://akaminas.github.io/vellum/}};
\end{tikzpicture}
\end{frame}
\end{document}
